\section{Discussion}
\label{sec:discussion}

\paragraph{The scaffold and the beam.} We started from an architecture
modeled on the mind (three networks, incubation, re-encounter), and
measurement returned something simpler and stronger than the architecture.
The DREAM scaffold works, but when it is taken apart almost all of its effect
lives in two operations that were not the noble ones: not eating one's literal
past, and being made, every few hundred tokens, to start a new sentence
somewhere else while remembering everything. We call this minimal recipe the
\emph{interrupted loop}. The elaborate parts either add
nothing (kicks, escalation), or cost something (forgetting costs connection),
or work opposite to their design (a literal return to the premise closes the
loop; salience is a good reader of the stream and a bad metronome for it).
Only because the whole scaffold was built and then dismantled do we know which
beam carries the weight, and which pieces are conditional: forgetting earns
its keep on a generator whose well is web boilerplate, and not on one that
drifts in prose.

\paragraph{Surface novelty and idea novelty are different quantities.} The
anti-probable sampler doubles $n$-gram novelty against the training corpus and
does nothing to judged surprise or connection; the interruption does nothing
special to $n$-grams and lifts judged surprise from 0.3 to 3--6. Inside the
network the two live at different depths: what the judge calls surprise is
lexical departure over an intact deep context, not a deep representational
shift, and the interruption that works barely moves the deep state. This is
also why the neat expectation ``novelty = distance from what came before''
fails: within interrupted streams, deep departure from the stream's own past
is \emph{negatively} related to judged surprise. The reader rewards the window
that is new on the surface and continuous underneath, which is what a good
turn in a story is.

\paragraph{Rhythm.} The frequency result is, we think, the most usable one:
a stream cut every 75 tokens develops nothing; the same break that scores 3.1
after a 150-token thread scores 5.3--5.9 after a 300--900-token thread; and
the yield decays over the next few hundred tokens. Surprise needs a developed
background to break, and interruption is a resource with a rate. The rhythm
we found, develop for a few hundred tokens and then break away, is the
incubation structure in miniature, the coffee break of a working mind: leave,
and be made to come back on your own.

\paragraph{What this is not.} It is not a claim that these streams are good
literature; most windows are judged well below the midpoint, and the texts are
what an unsupervised base model produces. It is not a claim about scale; the
generators are 8--30B parameters at 8 bits. It is a claim about \emph{where}
the novelty a base model can produce with no task comes from, made with
controls, three generator families and two judge families, and it points at
the loop's rhythm rather than at the two places where effort is usually
spent.
